\selectlanguage{ngerman}%
\frenchspacing%
\begin{ledgroupsized}[r]{120mm}%
\footnotesize%
\pstart%
\noindent%
\textbf{Überlieferung:}%
\pend%
\end{ledgroupsized}%
\begin{ledgroupsized}[r]{114mm}%
\footnotesize%
\pstart%
\parindent -6mm%
\makebox[6mm][l]{\textit{L}}%
Aufzeichnung:
LH XXXVII 4 Bl.\;162--163.
Ein Bogen 4\textsuperscript{o};
Ränder beschnitten;
Papiererhaltungsmaßnahmen.
Drei ganzseitig beschriebene Seiten;
Bl.\;163\,v\textsuperscript{o} unbeschrieben.
Der Text von N.\;\ref{57041} setzt die Aufzeichnung N.\;\ref{57040} unmittelbar fort.
\pend%
\end{ledgroupsized}%
\selectlanguage{latin}%
\frenchspacing%
\vspace*{5mm}%
\normalsize%
\pstart%
\noindent%
\lbrack%
\href{\Scan{LH_37_04_162r}}%
{162\,r\textsuperscript{o}}%
\rbrack%
\hfill%
\edtext{M.\;P.%
\indexS{motus perpetuus}%
}{%
\lemma{M.\;P.\,}%
\Cfootnote{\textit{motus perpetuus}
(so auch im Folgenden).%
}}
per Thermom.%
\indexS{thermometrum}
Barometr.%
\indexS{barometrum}
\edtext{%
Hygrometr.%
\indexS{hygrometrum}%
\hfill%
\phantom{\lbrack162\,r\textsuperscript{o}\rbrack}%
\\
\phantom{\lbrack162\,r\textsuperscript{o}\rbrack}%
\hfill%
Scheda 2%
\indexS{scheda}%
}{%
\lemma{Hygrometr.}%
\Bfootnote{%
\textit{(1)}~Pars %
\textit{(2)}~Scheda 2%
~\textit{L}%
}}%
\hfill%
\phantom{\lbrack162\,r\textsuperscript{o}\rbrack}%
\pend%
\vspace{0.5em}%
\pstart%
\noindent%
Pro
\edtext{%
Motore imperceptibili,%
\indexS{motor imperceptibilis}%
}{%
\lemma{Motore}%
\Bfootnote{%
\textit{(1)}~insensili %
\textit{(2)}~imperceptibili,%
~\textit{L}%
}}
perpetuo,%
\indexS{motor perpetuus}
forti,
nihil commodius reperio
quam aerem compressum%
\indexS{aer compressus}
\edtext{%
fortiter,
in spatio exiguo,
et postea superveniente calore%
\indexS{calor superveniens}%
}{%
\lemma{fortiter,}%
\Bfootnote{%
\textit{(1)}~et calore %
\textit{(2)}~in spatio %
\lbrack...\rbrack\ superveniente calore%
~\textit{L}%
}}
adhuc dilatatum.%
\indexS{aer dilatatus calore}
Nam
si aer communis%
\indexS{aer communis}
calore rarefactus%
\indexS{aer rarefactus calore}%
\indexS{calor superveniens}
spatium novum quaerit,%
\indexS{aer quaerens spatium novum}
quid faciet compressus%
\indexS{aer compressus}%
\lbrack?\rbrack\
Nimirum ita arbitror
\edtext{%
aerem%
\indexS{aer suscipiens vim caloris proportione densitatis}
vim ejusdem caloris%
\indexS{vis caloris}%
}{%
\lemma{aerem}%
\Bfootnote{%
\textit{(1)}~eundem %
\textit{(2)}~vim ejusdem caloris%
~\textit{L}%
}}
proportione suae densitatis%
\indexS{densitas aeris}
suscipere.%
\indexS{aer suscipiens vim caloris proportione densitatis}
Ita si
\edtext{%
in Tubo%
\indexS{tubus}%
}{%
\lemma{in}%
\Bfootnote{%
\textit{(1)}~pila\indexS{pila} %
\textit{(2)}~Tubo %
\textit{(3)}~vase\indexS{vas} %
\textit{(4)}~Tubo%
~\textit{L}%
}}
sit
\edtext{%
aer maxime rarefactus,%
\indexS{aer rarefactus}%
}{%
\lemma{aer}%
\Bfootnote{%
\textit{(1)}~rarefactus, aut ad sum %
\textit{(2)}~maxime rarefactus,%
~\textit{L}%
}}
quemadmodum contingit
cum antlia%
\indexS{antlia}
extrahitur,%
\indexS{aer extractus antlia}
\edtext{%
videmus calore externo%
\indexS{calor externus}
adhuc amplius rarefactum,%
\indexS{aer rarefactus}
vix tamen sensibilem vim exercere,%
\indexS{vis sensibilis}
alioqui omne barometrum%
\indexS{barometrum}
thermometri visibilis%
\indexS{thermometrum}
loco foret,
quod non fit,%
}{%
\lemma{videmus}%
\Bfootnote{%
\bars{tamen \textit{streicht Hrsg.}}
\textit{(1)}~ipsum %
\textit{(2)}~calore %
\textit{(a)}~adu %
\textit{(b)}~externo adhuc amplius rarefactum, %
\textit{(aa)}~vim tamen thermometri non habere, ac vix sensibiliter %
\textit{(bb)}~\textlangle mota\textrangle\ %
\textit{(cc)}~vix tamen sensibilem %
\textit{(aaa)}~vix %
\textit{(bbb)}~vim exercere, %
\lbrack...\rbrack\ non fit,%
~\textit{L}%
}}
si aere probe sit purgatum,%
\indexS{aer purgatus}
id est,
si in eo residuus aer%
\indexS{aer residuus}
sit rarissimus,%
\indexS{aer rarus}
qui hic pro nullo habetur.
Unde ita pronuntiandum arbitror
aerem calore dilatatum,%
\indexS{aer dilatatus calore}
tantam vim habere,%
\indexS{vis aeris dilatati}
quantam haberet,
\edtext{%
si hic se dilatandi%
}{%
\lemma{si}%
\Bfootnote{%
\textit{(1)}~haec di %
\textit{(2)}~hic se dilatandi%
~\textit{L}%
}}
conatus%
\indexS{conatus aeris dilatandi se}
a proprio aeris elastro,%
\indexS{elastrum aeris}
compressione scilicet antea facta,%
\indexS{compressio aeris}
veniret.
Quod ita mechanice%
\indexS{explicatio mechanica}
\edtext{%
explico%
\lbrack:\rbrack\
fingamus aerem%
\indexS{aer similis lanae}
esse lanae%
\indexS{lana}
similem,%
}{%
\lemma{explico,}%
\Bfootnote{%
\textit{(1)}~ponamus %
\textit{(2)}~fingamus aerem esse %
\textit{(1)}~lanam %
\textit{(2)}~lanae similem,%
~\textit{L}%
}}
calore autem%
\indexS{calor extendens filum}
unumquodque ejus filum se extendere atque explicare%
\indexS{filum extensum calore}
(ut lana humore),%
\indexS{lana extensa humore}%
\indexS{humor extendens lanam}
patet
quo compressior est lana,%
\indexS{lana compressa}%
\indexS{compressio lanae}
hoc magis humore dilatari%
\indexS{lana dilatata humore}
quia quo compressior est,
hoc plura in eodem loco sunt fila,%
\indexS{filum suscipiens dilatationem}%
\indexS{filum lanae}
quorum unumquodque dilatationem%
\indexS{dilatatio filorum lanae}
suscipit%
\lbrack,\rbrack\
eodem modo quo compressior est aer,%
\indexS{aer compressus}
hoc fortius se rarefactione%
\indexS{rarefactio aeris}%
\indexS{aer rarefactus}
explicabit.%
\indexS{aer explicatus rarefactione}
Et quoniam minimum decuplum
aeris%
\indexS{aer intrusus}
prioris intrudi%
\indexS{aer intrusus}
potest,
patet et vim decuplicari,%
\indexS{vis aeris decuplicata}
vis
\edtext{%
autem aeris ordinarii%
\indexS{aer ordinarius}
calore rarefacti%
\indexS{aer rarefactus calore}
est tanta
ut pondus%
\indexS{pondus aeris}
aeris incumbentis%
\indexS{aer incumbens}
elevet,%
}{%
\lemma{autem}%
\Bfootnote{%
\textit{(1)}~est tanta quae pondus aeris imminentis elevat %
\textit{(2)}~aeris ordinarii calore %
\textit{(a)}~rarefaciendi %
\textit{(b)}~rarefacti est %
\lbrack...\rbrack\ incumbentis elevet,%
~\textit{L}%
}}
id est,
si unius pedis%
\indexS{pes}
sit diameter follis,%
\indexS{follis ahibitus barometri}%
\indexS{follis thermometri}%
\indexS{diameter follis}
librarum%
\indexS{libra}
circiter~12.
Adeoque rarefacto aere%
\indexS{aer rarefactus calore}
saltem ita,
\edtext{%
ut prioris quintuplum%
}{%
\lemma{ut}%
\Bfootnote{%
\textit{(1)}~quadruplum prio %
\textit{(2)}~prioris quintuplum%
~\textit{L}%
}}
insit,
erit vis%
\indexS{vis aeris rarefacti}
librarum sexaginta,%
\indexS{libra}
quae objecto
\edtext{%
fortissimo rigidissimoque Elastro ferreo%
\indexS{elastrum ferreum}%
\indexS{elastrum rigidum}%
}{%
\lemma{fortissimo}%
\Bfootnote{%
\textit{(1)}~Elaterio\indexS{elaterium ferreum} %
\textit{(2)}~rigidissimoque Elastro ferreo%
~\textit{L}%
}}
moderanda est.%
\indexS{vis aeris moderanda elastro}
Sed ponamus esse multo minorem,
nihil refert,
quia etiam aucta compressione compensabitur%
\indexS{compressio aeris}
quod calculo deest%
\lbrack;\rbrack\
\indexS{calculus}
sed etsi non nisi ordinarius,
id
\edtext{%
est 12 librarum,%
}{%
\lemma{est}%
\Bfootnote{%
\textit{(1)}~60 lib %
\textit{(2)}~12 librarum,%
~\textit{L}%
}}
esset aer,%
\indexS{aer ordinarius}
nec compressio%
\indexS{compressio aeris}
adderet quicquam,
nihilominus foret
duplicata diametro%
\indexS{diameter follis}
librarum 48 seu prope 50,%
\indexS{libra}
si scilicet diameter sit 2 pollicum.%
\indexS{pollex}
Quae omnia plus quam abunde sufficere
\edtext{%
videntur,
\lbrack%
\href{\Scan{LH_37_04_162v}}%
{162\,v\textsuperscript{o}}%
\rbrack\
praesertim%
}{%
\lemma{videntur,}%
\Bfootnote{%
\textit{(1)}~praesertim %
\textit{(2)}~tanta autem %
\textit{(3)}~\lbrack\href{\Scan{LH_37_04_162v}}{162\,v\textsuperscript{o}}\rbrack\ praesertim%
~\textit{L}%
}}
cum mutatio circa calorem%
\indexS{mutatio circa calorem et frigus}%
\indexS{calor aeris}
et frigus%
\indexS{frigus aeris}
quotidie obtingat.
Hinc
tametsi exiguum sit spatium
in quo sit motus,
magnitudo tamen ponderis%
\indexS{magnitudo ponderis}%
\indexS{pondus elevandum}
hunc defectum supplet,
nam 12 libras%
\indexS{libra}
\edtext{%
per dimidii pollicis%
\indexS{pollex}
altitudinem%
}{%
\lemma{per}%
\Bfootnote{%
\textit{(1)}~dimidium pollicem, %
\textit{(a)}~de %
\textit{(b)}~qui t %
\textit{(2)}~dimidii pollicis altitudinem%
~\textit{L}%
}}
aut semilibram%
\indexS{semilibra}
per pedis%
\indexS{pes}
altitudinem descendere
perinde
\edtext{%
est.
At semilibra
per pedis altitudinem%
}{%
\lemma{est.}%
\Bfootnote{%
\textit{(1)}~Librae autem %
\textit{(2)}~At semilibra per %
\textit{(a)}~pollicis %
\textit{(b)}~pedis altitudinem%
~\textit{L}%
}}
descendente
horologium aliquot diebus duraturum%
\indexS{horologium duraturum aliquot diebus}
non difficulter credo fiet.%
\pend%
\pstart%
Caeterum
\edtext{%
si motum fortem
in usus domesticos%
\indexS{motus fortis}%
}{%
\lemma{si}%
\Bfootnote{%
\textit{(1)}~in magno %
\textit{(2)}~motum %
\textit{(a)}~fortiorem %
\textit{(b)}~fortem in usus domesticos%
~\textit{L}%
}}
et Molendini speciem quaeramus%
\indexS{species molendini quaesita}%
\indexS{molendinum quaesitum}
post
\edtext{%
ventaneum%
\indexS{molendinum ventaneum perpetuum}
\edtext{%
perpetuum priore scheda%
}{%
\lemma{perpetuum}%
\Bfootnote{%
\textit{(1)}~alias %
\textit{(2)}~priore scheda%
~\textit{L}%
}}
descriptum,%
}{%
\lemma{ventaneum \lbrack...\rbrack\ descriptum}%
\Cfootnote{%
Siehe N.\;\ref{57040}, \refpassage{LH_37_04_160-161_160r_molvent-1}{LH_37_04_160-161_161r_molvent-2}.%
}}
hoc
quod calore animatur,%
\indexS{molendinum animatum calore}%
\indexS{calor animans molendinum}
fortissimum arbitror.
\edtext{%
Esto enim follis%
\indexS{follis ad circumagendum molendinum}
latitudinis pedalis,
aer ordinarius%
\indexS{aer ordinarius}
qui intus est,%
\indexS{aer in folle}
sustinet aeris incumbentis%
\indexS{aer incumbens}
libras 1700.%
\indexS{libra}%
}{%
\lemma{Esto \lbrack...\rbrack\ 1700}%
\Cfootnote{%
Siehe N.\;\ref{57040}, \refpassage{LH_37_04_160-161_161r_pasca-1}{LH_37_04_160-161_161r_pasca-2}.%
}}
Sit decuplo compressior,%
\indexS{aer compressus}
continebitur vi%
\indexS{vis aeris compressi}
librarum 17000,%
\indexS{libra}
id est,
aer incumbens%
\indexS{aer incumbens}
libras 1700 facit
et pondera adjecta reliqua%
\edtext{}{%
\lemma{\textit{Zwischen den Zeilen, umkreist:}}%
\Afootnote{%
\begin{tabular}[t]{r}%
$17000 \ -$%
\\%
$1700\phantom{\ -}$%
\\%
\hline%
$15300\phantom{\ -}$%
\end{tabular}%
}}
15300 libras facient.%
\indexS{libra}
Ponamus jam calore dilatari%
\indexS{aer dilatatus calore}%
\indexS{calor dilatans aerem}
qui talis sit
ut spatium occupare conetur inclusus aer,%
\indexS{aer inclusus folli}%
\indexS{aer conans occupare spatium}
quod sit ad prius
\edtext{%
ut 6 ad 5.%
}{%
\lemma{ut}%
\Bfootnote{%
\textit{(1)}~12 ad 9 %
\textit{(2)}~6 ad 5.%
~\textit{L}%
}}
Ponatur et altitudo follis%
\indexS{follis ad circumagendum molendinum}
\edtext{%
esse 5 pedum,%
}{%
\lemma{esse}%
\Bfootnote{%
\textit{(1)}~9 pedum, utique %
\textit{(2)}~5 pedum,%
~\textit{L}%
}}
(latitudinem posuimus unius pedis,%
\indexS{pes}
forma est cylindrica)%
\indexS{forma follis cylindrica}
patet %
nunc occupaturum spatium sex pedum,%
\indexS{pes}
atque ita perinde fore,
ac si 17000 librae%
\indexS{libra}
per unius pedis%
\indexS{pes}
spatium descenderent,
quae vis%
\indexS{vis aeris dilatati}
quam immanis sit%
\indexS{vis immanis}
manifestum
\edtext{%
est.
Cessante rursus calore,%
\indexS{calor cessans}%
}{%
\lemma{est}%
\Bfootnote{%
\textit{(1)}~, haec porro mutatio %
\textit{(2)}~. Cessante rursus calore,%
~\textit{L}%
}}
pondera illa maxima appensa%
\indexS{pondus appensum molendino}
suopte nisu%
\indexS{nisus ponderis appensi}
follem in priorem angustiam restituent,%
\indexS{follis restitutus}
et hac reciprocatione fiet,%
\indexS{reciprocatio}
ut semper motum maximum in promtu habeamus.%
\indexS{motus in promptu}
\pend%
\pstart%
Hinc
\edtext{%
etiam mira%
}{%
\lemma{etiam}%
\Bfootnote{%
\textit{(1)}~mirum %
\textit{(2)}~mira%
~\textit{L}%
}}
vis patet%
\indexS{vis mira}
simplicis luminis candelae,%
\indexS{vis luminis candelae}%
\indexS{lumen candelae}%
\indexS{candela}
imo etiam manus calidae,%
\indexS{vis manus calidae}%
\indexS{manus calidae}
haec enim imposita
\edtext{%
thermometro aperto%
\indexS{thermometrum apertum}
facit aerem%
}{%
\lemma{thermometro}%
\Bfootnote{%
\textit{(1)}~facit aer %
\textit{(2)}~aperto facit aerem%
~\textit{L}%
}}
occupare spatium majus.%
\indexS{aer occupans spatium majus}
Perinde ergo est,
ac si aer id propria compressione%
\indexS{compressio aeris}
conaretur%
\lbrack;\rbrack\
quanta autem aeris compressi%
\indexS{aer compressus}
vis sit,%
\indexS{vis aeris compressi}
patet%
\lbrack:\rbrack\
nempe cum ordinarii%
\indexS{aer ordinarius}
Elastrum%
\indexS{elastrum aeris}
vim habeat tantam%
\indexS{vis elastri aeris}
quanta est aeris incumbentis pondus,%
\indexS{pondus aeris incumbentis}%
\indexS{aer incumbens}
patet compressum%
\indexS{aer compressus}
cum mole%
\indexS{moles aeris compressi}
et vim augere.%
\indexS{vis aeris compressi}
Hinc minus mirabimur,
tantam vim esse glaciationis,%
\indexS{vis glaciationis}%
\indexS{glaciatio}
ut fortissima vasa disrumpat,%
\indexS{vas disruptum}
cum tanta sit vis luminis candelae.%
\indexS{vis luminis candelae}%
\indexS{lumen candelae}%
\indexS{candela}
Inquirendum
an calor distendat aerem,%
\indexS{calor distendens aerem}%
\indexS{aer distensus calore}
an vero ipsum velut congelatum%
\indexS{aer congelatus}
quasi fundens,%
\indexS{calor fundens aerem}
potestatem det partibus aeris%
\indexS{potestas aeris explicandi se}%
\indexS{aer explicans se}
se explicandi Elastro spontaneo.%
\indexS{elastrum aeris spontaneum}
Prius verisimilius.
\lbrack%
\href{\Scan{LH_37_04_163r}}%
{163\,r\textsuperscript{o}}%
\rbrack%
\pend%
\pstart%
Ope porro
\edtext{%
follis mei chalybei exigui,%
\indexS{follis chalybeus}%
}{%
\lemma{follis \lbrack...\rbrack\ exigui}%
\Cfootnote{%
Siehe N.\;\ref{57040}, \refpassage{LH_37_04_160-161_161v_folfer-1}{LH_37_04_160-161_161v_folfer-2}.%
}}
(musculorum fabricam referentis)%
\indexS{fabrica musculorum}%
\indexS{musculus}
non tantum M.\;P.%
\indexS{motus perpetuus}
in horologio portatili%
\indexS{horologium portatile}
exhibebo,
sed et pulcherrimum dabo aerometrum,%
\indexS{aerometrum sensibile}
quod in exiguo non minus sensile sit,
quam
\edtext{%
\edlabel{LH_37_04_162-163_163r_hooke-1}%
majus illud Anglicano more elaboratum.%
\indexS{aerometrum Anglicanum}%
\edlabel{LH_37_04_162-163_163r_hooke-2}%
}{%
\lemma{majus \lbrack...\rbrack\ elaboratum}%
\Cfootnote{%
Anspielung auf das von Hooke erfundene Radbarometer.
Siehe \textsc{R.\;Hooke},
\textit{Micrographia}, Preface
(London 1665, S.\:[c2]--[c3] u.\;Tafel I, Abb.\;I);%
\cite{00061}
\textsc{Ders.},
\glqq A new Contrivance of Wheel-Barometer\grqq\
(\textit{PT}, I, 4.\;Juni 1666,
S.\;218\,f. u.\;Abb.\;1 vor S.\;215).%
\cite{01523}
Auf das \glqq englische\grqq\ Radbarometer beziehen sich in diesem Band auch die Aufzeichnung
N.\;\ref{57038}
(\refpassage{LH_37_04_157-158_157v_udz-1}{LH_37_04_157-158_157v_udz-2})
vom November 1677
sowie die spätere Aufzeichnung
N.\;\ref{53451}
(\refpassage{LH_37_04_171_171r_hooke-1}{LH_37_04_171_171r_hooke-2}).%
}}
Scilicet
\edtext{%
sit follis iste%
\indexS{follis chalybeus}
latitudine pollicis,%
\indexS{pollex}
erit pondus incumbens%
\indexS{pondus incumbens folli}
librarum 12,%
\indexS{libra}%
}{%
\lemma{sit \lbrack...\rbrack\ librarum~12}%
\Cfootnote{%
Siehe N.\;\ref{57040}, \refpassage{LH_37_04_160-161_161v_sgi-1}{LH_37_04_160-161_161v_sgi-2}.%
}}
respondens 30 pollicibus argenti vivi,%
\indexS{argentum vivum}
hujus pars
\edtext{%
duodecima vel%
}{%
\lemma{duodecima vel}%
\Bfootnote{%
\textit{erg.\;L}%
}}
decima fere subjecta mutationi,%
\indexS{mutatio argenti vivi}
id est pollices duo vel tres,%
\indexS{pollex}
restabit ergo libra una;%
\indexS{libra}
\edtext{%
\edtext{%
sed
\edlabel{LH_37_04_162-163_163r_hooke-3}%
in%
}{%
\lemma{sed}%
\Bfootnote{%
\textit{(1)}~cum %
\textit{(2)}~in%
~\textit{L}%
}}
Anglicano illo%
\indexS{aerometrum Anglicanum}
hi pollices duo vel tres rursus
in 200 partes sunt divisi,%
\edlabel{LH_37_04_162-163_163r_hooke-4}%
}{%
\lemma{sed \lbrack...\rbrack\ divisi}%
\Cfootnote{%
Vgl.
\textsc{Hooke},
\textit{Micrographia}, Preface
(S.\:[c3]).%
\cite{00061}
}}
\edtext{%
itaque ducentesima librae pars%
\indexS{libra}
indicis gyrationi%
\indexS{gyratio indicis aerometri}%
\indexS{index aerometri}
sufficere debet,%
}{%
\lemma{itaque}%
\Bfootnote{%
\textit{(1)}~vigesima librae pars %
\textit{(2)}~ducentesima librae pars %
\textit{(a)}~tantum %
\textit{(b)}~tantum aget %
\textit{(c)}~acti %
\textit{(d)}~indicis gyrationi sufficere debet,%
~\textit{L}%
}}
ut non minus sensibile habeamus%
\indexS{aerometrum sensibile}
quam est anglicanum,%
\indexS{aerometrum Anglicanum}%
\edtext{}{%
\lemma{\textit{Hilfsrechnung zum Haupttext:}}%
\Afootnote{%
$\left.\displaystyle\frac{16}{200}\right|\displaystyle\frac{2}{25}$%
\ \textsuperscript{\lbrack a\rbrack}%
\\%
[1ex]%
{\footnotesize%
\textsuperscript{\lbrack a\rbrack}~%
\textit{(1)}\:$\left.\displaystyle\frac{12}{200}\right|\displaystyle\frac{3}{50}$ %
\textit{(2)}\:$\left.\displaystyle\frac{16}{200}\right|\displaystyle\frac{2}{25}$%
~\textit{L}%
}%
}}
id est pars circiter tertia unciae,%
\indexS{uncia}
quarum 16 faciunt libram,%
\indexS{libra}
quod pondus
circumagendo indici levissimo%
\indexS{index aerometri}
sufficere arbitror,%
\indexS{pondus sufficiens ad indicem circumagendum}
cum nulla alia vi sit opus.
Aliud a pondere aeris%
\indexS{pondus aeris}
ducitur ad domesticos usus%
\indexS{usus domesticus}
mire validum,%
\indexS{aerometrum validum ad usus domesticos}
scilicet follis vacuus distentus,%
\indexS{follis vacuus distentus}
qui%
\lbrack,\rbrack\
cum opus
relaxatus%
\lbrack,\rbrack\
maximam vim exercebit.
\pend%
\pstart%
Si tanta
\edtext{%
vis%
\indexS{vis inclusa moli aeris}%
\indexS{vis pendens a folle thermometrico}
quam%
}{%
\lemma{vis}%
\Bfootnote{%
\textit{(1)}~tam %
\textit{(2)}~quam%
~\textit{L}%
}}
\edtext{%
a folle Thermometrico%
\indexS{follis thermometricus}
pendere%
}{%
\lemma{quam \lbrack...\rbrack\ ostendi}%
\Cfootnote{%
Siehe N.\;\ref{57040}, \refpassage{LH_37_04_160-161_161v_fltherm}{LH_37_04_160-161_161v_fltherm}.%
}}
\edtext{%
ostendi,
tam exiguae moli%
\indexS{moles aeris compressi}%
}{%
\lemma{ostendi,}%
\Bfootnote{%
\textit{(1)}~aerii %
\textit{(2)}~tam %
\textit{(a)}~levi %
\textit{(b)}~exiguae moli%
~\textit{L}%
}}
includi potest,%
\indexS{vis inclusa moli aeris}
non desperem fieri posse machinam%
\indexS{machina movens se}
se moventem per spatium
\edtext{%
et tempus%
}{%
\lemma{et tempus}%
\Bfootnote{%
\textit{erg.\;L}%
}}
satis longum.
Quis enim non videt
libras 17000%
\indexS{libra}
per unum pedem%
\indexS{pes}
descendentes
\edtext{%
facile pondus 300 librarum%
\indexS{pondus hominis}
quale est homo,%
\indexS{homo}%
}{%
\lemma{facile}%
\Bfootnote{%
\textit{(1)}~homi %
\textit{(2)}~pondus %
\textit{(a)}~400 %
\textit{(b)}~300 librarum quale est homo,%
~\textit{L}%
}}
tracturas per
\rule[-3.0mm]{0pt}{8.4mm}%
$\displaystyle\frac{170}{3}$
pedis,%
\indexS{pes}
seu per pedes 60 circiter.
Si follis sit tripedalis latitudinis,
sustinebit libras 15300
($9\smallfrown 1700$),
et si aer sit decuplo compressior solito,%
\indexS{aer compressus}
erit intus vis librarum 153000,
quae descendens per unum pedem
pondus 1000 librarum%
\indexS{libra}
\edtext{%
ducet facile per%
}{%
\lemma{ducet}%
\Bfootnote{%
\textit{(1)}~per %
\textit{(2)}~facile per%
~\textit{L}%
}}
pedes 150.%
\indexS{pes}
Sin descendat per tres pedes
erit vis ut
$153000\smallfrown 3$
seu 459000 librarum descendentium per unum
\edtext{%
pedem,
et facile trahet pondus 1000 librarum%
\indexS{libra}
per pedes 450.%
\indexS{pes}
Difficultas%
\indexS{difficultas}%
}{%
\lemma{pedem,}%
\Bfootnote{%
\textit{(1)}~seu 459 %
\textit{(2)}~et facile trahet pondus 1000 %
\textit{(a)}~pedum per libras %
\textit{(b)}~p %
\textit{(c)}~librarum per pedes 450. %
\textit{(aa)}~Res %
\textit{(bb)}~Difficultas%
~\textit{L}%
}}
eo redit
ut munitae satis ad vim sustinendam%
\indexS{vis sustinenda}
sint machinae,%
\indexS{machina munita ad vim sustinendam}
quod tamen in exiguis portatilibus,
maxime admirandis,
minime metuendum.
\pend%
\selectlanguage{latin}%
\frenchspacing%
\normalsize%
